\question[3]
Erläutere das Problem des Routings zwischen Netzen:
Warum reicht es nicht, dass jedes Gerät nur seine eigene lokale Netzadresse kennt,
wenn Daten an ein Gerät in einem völlig anderen Netz (z.\,B. auf einem anderen Kontinent)
gesendet werden sollen?

\begin{solutionbox}{6cm}

Ein Gerät kennt nur die Geräte in seinem eigenen lokalen Netz
direkt. Liegt das Ziel in einem anderen Netz, muss das
Datenpaket über mehrere Zwischenstationen (Router) transportiert
werden, die jeweils nur einen Ausschnitt des Gesamtnetzes
kennen.

Das Problem des Routings besteht darin, für jedes Paket auf
jedem beteiligten Router zu entscheiden, über welche
Verbindung (welchen nächsten Knoten) das Paket weitergeleitet
werden muss, damit es auf einem möglichst guten Weg
schließlich sein Ziel erreicht - obwohl kein einzelner Router
den kompletten Weg von Anfang bis Ende kennt.

\end{solutionbox}